\documentclass[11pt,a4paper]{article}
\usepackage[utf8]{inputenc}
\usepackage[margin=0.8in]{geometry}
\usepackage{amsmath,amssymb}
\usepackage{graphicx}
\usepackage{xcolor}
\usepackage{booktabs}
\usepackage{listings}
\usepackage{fancyhdr}
\usepackage{hyperref}

\pagestyle{fancy}
\fancyhf{}
\rhead{\textbf{DSP Practical Lab Manual}}
\lhead{Experiments 1 \& 2: Theory, Formulas \& Derivations}
\rfoot{Page \thepage}

\definecolor{codegreen}{rgb}{0,0.6,0}
\definecolor{codegray}{rgb}{0.5,0.5,0.5}
\definecolor{codepurple}{rgb}{0.58,0,0.82}
\definecolor{backcolour}{rgb}{0.96,0.97,0.98}
\definecolor{headcolor}{rgb}{0.1,0.25,0.5}

\lstdefinestyle{scilabstyle}{
    backgroundcolor=\color{backcolour},   
    commentstyle=\color{codegreen}\itshape,
    keywordstyle=\color{headcolor}\bfseries,
    numberstyle=\tiny\color{codegray},
    stringstyle=\color{codepurple},
    basicstyle=\ttfamily\small,
    breakatwhitespace=false,         
    breaklines=true,                 
    captionpos=b,                    
    keepspaces=true,                 
    numbers=left,                    
    numbersep=5pt,                  
    showspaces=false,                
    showstringspaces=false,
    showtabs=false,                  
    tabsize=2,
    frame=single,
    rulecolor=\color{codegray}
}

\lstset{style=scilabstyle}

\begin{document}

\begin{center}
    {\Large \textbf{\textcolor{headcolor}{DIGITAL SIGNAL PROCESSING (DSP) LABORATORY}}}\\[0.2cm]
    {\large \textbf{Comprehensive Practical Guide: Experiments 1 \& 2}}\\[0.1cm]
    \rule{\linewidth}{0.5mm}
\end{center}

\vspace{0.2cm}

% ==============================================================================
% EXPERIMENT 1
% ==============================================================================
\section*{\textcolor{headcolor}{EXPERIMENT 1: Generation of Elementary Discrete-Time Sequences}}

\subsection*{1. Aim}
To generate and plot elementary discrete-time sequences including the Unit Impulse sequence, Unit Step sequence, Unit Ramp sequence, and Real-Valued Exponential sequence using Scilab.

\subsection*{2. Theory \& Mathematical Definitions}

Discrete-time signals are represented mathematically as sequences of numbers indexed by an integer $n \in \mathbb{Z}$. The elementary signals serve as foundational building blocks for analyzing linear discrete-time systems.

\subsubsection*{A. Unit Sample / Impulse Sequence $\delta[n]$}
The unit impulse sequence (Kronecker delta) is defined as:
\begin{equation}
    \delta[n] = 
    \begin{cases} 
        1, & n = 0 \\ 
        0, & n \ne 0 
    \end{cases}
\end{equation}
\textbf{Fundamental Property (Sifting Property):} Any arbitrary sequence $x[n]$ can be represented as a linear combination of shifted impulses:
\begin{equation}
    x[n] = \sum_{k=-\infty}^{\infty} x[k]\,\delta[n - k]
\end{equation}

\subsubsection*{B. Unit Step Sequence $u[n]$}
The unit step sequence represents a signal that switches on at $n=0$ with unit amplitude:
\begin{equation}
    u[n] = 
    \begin{cases} 
        1, & n \ge 0 \\ 
        0, & n < 0 
    \end{cases}
\end{equation}
\textbf{Relationship with Impulse Sequence:}
\begin{equation}
    u[n] = \sum_{k=-\infty}^{n} \delta[k] \quad \text{and} \quad \delta[n] = u[n] - u[n-1]
\end{equation}

\subsubsection*{C. Unit Ramp Sequence $r[n]$}
The unit ramp sequence increases linearly with index $n$ for non-negative indices:
\begin{equation}
    r[n] = 
    \begin{cases} 
        n, & n \ge 0 \\ 
        0, & n < 0 
    \end{cases} 
    \implies r[n] = n \cdot u[n]
\end{equation}
\textbf{First Difference Property:} $r[n] - r[n-1] = u[n-1]$.

\subsubsection*{D. Real-Valued Exponential Sequence $x[n] = a^n$}
The causal real-valued exponential sequence is defined as:
\begin{equation}
    x[n] = a^n \cdot u[n]
\end{equation}
Depending on the base parameter $a \in \mathbb{R}$:
\begin{itemize}
    \item $0 < a < 1$ (e.g., $a = 0.8$): \textbf{Monotonically decaying sequence} ($\lim_{n\to\infty} a^n = 0$).
    \item $a > 1$: \textbf{Monotonically growing sequence} ($\lim_{n\to\infty} a^n = \infty$).
    \item $-1 < a < 0$: \textbf{Alternating decaying sequence}.
    \item $a < -1$: \textbf{Alternating growing sequence}.
\end{itemize}

\subsection*{3. Scilab Implementation Code}
\begin{lstlisting}[language=Scilab]
// Experiment 1: Elementary Discrete-Time Signals
clear;
clc;

n = -10:50;  // Range of sample indices

// 1. Unit Sample Sequence delta[n]
x0 = zeros(1, length(n));
x0(n == 0) = 1;
subplot(2, 2, 1);
plot2d3(n, x0);
title("Unit Sample Sequence \delta[n]");
xlabel("Sample index n"); ylabel("\delta[n]");

// 2. Unit Step Sequence u[n]
x1 = ones(1, length(n));
x1(n < 0) = 0;
subplot(2, 2, 2);
plot2d3(n, x1);
title("Unit Step Sequence u[n]");
xlabel("Sample index n"); ylabel("u[n]");

// 3. Ramp Sequence r[n]
x2 = n;
x2(n < 0) = 0;
subplot(2, 2, 3);
plot2d3(n, x2);
title("Ramp Sequence r[n]");
xlabel("Sample index n"); ylabel("r[n]");

// 4. Real Exponential Sequence x[n] = (0.8)^n
x3 = (0.8).^n;
x3(n < 0) = 0;
subplot(2, 2, 4);
plot2d3(n, x3);
title("Real Exponential Sequence (a=0.8)");
xlabel("Sample index n"); ylabel("x[n]");
\end{lstlisting}

\subsection*{4. Results \& Observations}
\begin{enumerate}
    \item $\delta[n]$ is non-zero only at $n=0$ with amplitude $1$.
    \item $u[n]$ remains constant at $1$ for all $n \ge 0$.
    \item $r[n]$ exhibits a unit slope for $n \ge 0$.
    \item $x[n] = (0.8)^n$ decays smoothly towards $0$ as $n$ increases, confirming stability.
\end{enumerate}

\newpage

% ==============================================================================
% EXPERIMENT 2
% ==============================================================================
\section*{\textcolor{headcolor}{EXPERIMENT 2: Linear Convolution of Discrete-Time Sequences}}

\subsection*{1. Aim}
To compute and plot the linear convolution of two discrete-time rectangular pulse sequences:
\begin{itemize}
    \item \textbf{Part A:} Using Scilab's built-in function \texttt{convol()}.
    \item \textbf{Part B:} Without using built-in functions, implementing the mathematical sum-of-products formula via nested for-loops.
\end{itemize}

\subsection*{2. Theory \& Mathematical Formulation}

Linear convolution describes the input-output relationship of any Linear Time-Invariant (LTI) system. If an LTI system has an impulse response $h[n]$ and is excited by an arbitrary input $x[n]$, the output $y[n]$ is the convolution sum:
\begin{equation}
    \mathbf{y[n] = x[n] * h[n] = \sum_{k=-\infty}^{\infty} x[k] \, h[n - k]}
\end{equation}

\subsubsection*{Fundamental Properties of Linear Convolution}
\begin{enumerate}
    \item \textbf{Commutative Property:} $x[n] * h[n] = h[n] * x[n]$
    \item \textbf{Associative Property:} $\big(x[n] * h_1[n]\big) * h_2[n] = x[n] * \big(h_1[n] * h_2[n]\big)$
    \item \textbf{Distributive Property:} $x[n] * \big(h_1[n] + h_2[n]\big) = x[n] * h_1[n] + x[n] * h_2[n]$
\end{enumerate}

\subsection*{3. Mathematical Derivations \& Calculations}

\paragraph{Step 1: Length and Sample Range Derivation}
Let sequence $x_0[n]$ have length $N_1$ spanning from index $n_1$ to $n_2$:
\begin{equation}
    N_1 = n_2 - n_1 + 1
\end{equation}
Let sequence $x_1[n]$ have length $N_2$ spanning from index $m_1$ to $m_2$:
\begin{equation}
    N_2 = m_2 - m_1 + 1
\end{equation}
Then the convolution output $y[n]$ spans from:
\begin{align}
    \mathbf{n_{start}} &= n_1 + m_1 \\
    \mathbf{n_{end}}   &= n_2 + m_2
\end{align}
And the total number of output samples $N_y$ is strictly given by:
\begin{equation}
    \mathbf{N_y = N_1 + N_2 - 1}
\end{equation}

\paragraph{Step 2: Applied to the Laboratory Experiment Sequences}
In the laboratory experiment, both sequences are defined over index range $n = [-10, 10]$:
\begin{equation}
    x_0[n] = x_1[n] = 
    \begin{cases} 
        1, & -5 \le n \le 5 \quad (\text{pulse width } W = 11 \text{ samples}) \\ 
        0, & \text{otherwise in } -10 \le n \le 10 
    \end{cases}
\end{equation}
\begin{itemize}
    \item Length of input arrays: $N_1 = 10 - (-10) + 1 = 21$ samples.
    \item Length of impulse response array: $N_2 = 21$ samples.
    \item Output length: $N_y = 21 + 21 - 1 = \mathbf{41\text{ samples}}$.
    \item Output index range:
    \begin{align}
        n_{start} &= -10 + (-10) = \mathbf{-20} \nonumber \\
        n_{end}   &= 10 + 10 = \mathbf{+20} \nonumber \\
        \mathbf{ny} &= [-20, +20]
    \end{align}
\end{itemize}

\paragraph{Step 3: Analytical Derivation of Triangular Waveform Output}
The convolution of two identical rectangular pulses of width $W = 11$ samples yields a \textbf{symmetric triangular sequence}:
\begin{itemize}
    \item \textbf{Peak Amplitude at $n = 0$:}
    \begin{equation}
        y[0] = \sum_{k=-5}^{5} x_0[k]\,x_1[0 - k] = \sum_{k=-5}^{5} (1)(1) = \mathbf{11}
    \end{equation}
    \item \textbf{Analytical Formula for Output Values:}
    \begin{equation}
        y[n] = 
        \begin{cases} 
            11 - |n|, & |n| \le 10 \\ 
            0, & |n| > 10 
        \end{cases}
    \end{equation}
\end{itemize}

\paragraph{Step 4: Algorithmic Derivation for Part B (Nested For-Loops)}
In Scilab, array indices are 1-based ($1 \le i \le N_1$ and $1 \le j \le N_2$).
The output sample index mapping for the product $x_0(i) \cdot x_1(j)$ corresponds to output index:
\begin{equation}
    k = i + j - 1
\end{equation}
When $i=1$ and $j=1$, $k = 1 + 1 - 1 = 1$ (first sample). When $i=N_1$ and $j=N_2$, $k = N_1 + N_2 - 1 = N_y$ (last sample).
Thus, the inner accumulation step is:
\begin{equation}
    y(i + j - 1) = y(i + j - 1) + x_0(i) \cdot x_1(j)
\end{equation}

\subsection*{4. Scilab Implementation Code}

\subsubsection*{Part A: Using Built-in Function \texttt{convol()}}
\begin{lstlisting}[language=Scilab]
// Experiment 2 - Part A: Linear Convolution using convol()
clear;
clc;

n = -10:10;
x0 = zeros(1, length(n));
x0(-5 <= n & n <= 5) = 1;

x1 = zeros(1, length(n));
x1(-5 <= n & n <= 5) = 1;

y = convol(x0, x1);
ny = -20:20;  // Length = 21 + 21 - 1 = 41

subplot(3, 1, 1);
plot2d3(n, x0);
title("Input Pulse x0[n]");
xlabel("n"); ylabel("x0[n]");

subplot(3, 1, 2);
plot2d3(n, x1);
title("Input Pulse x1[n]");
xlabel("n"); ylabel("x1[n]");

subplot(3, 1, 3);
plot2d3(ny, y);
title("Linear Convolution Output y[n] = x0[n] * x1[n]");
xlabel("n"); ylabel("y[n]");
\end{lstlisting}

\subsubsection*{Part B: Manual Implementation using For-Loops}
\begin{lstlisting}[language=Scilab]
// Experiment 2 - Part B: Linear Convolution using Nested Loops
clear;
clc;

n = -10:10;
x0 = zeros(1, length(n));
x0(-5 <= n & n <= 5) = 1;

x1 = zeros(1, length(n));
x1(-5 <= n & n <= 5) = 1;

N1 = length(x0);
N2 = length(x1);
Ny = N1 + N2 - 1;

y = zeros(1, Ny);
for i = 1:N1
    for j = 1:N2
        y(i + j - 1) = y(i + j - 1) + x0(i) * x1(j);
    end
end

ny = -20:20;
plot2d3(ny, y);
title("Linear Convolution Output (Computed via For-Loops)");
xlabel("Sample index n"); ylabel("y[n]");
\end{lstlisting}

\subsection*{5. Results \& Discussion}
\begin{enumerate}
    \item Both Part A (\texttt{convol}) and Part B (nested loops) produce identical output arrays of length $41$ spanning from $n = -20$ to $n = +20$.
    \item The convolution of two rectangular pulses of width $11$ samples produces a symmetric triangular pulse peaking at $n = 0$ with amplitude $11$.
\end{enumerate}

\end{document}
